
\section{Giới thiệu}

% ===== SLIDE CHUYỂN TIẾP PHẦN 1 =====
\begin{frame}{}
  \centering
  \vfill
  {\Huge \textbf{Phần 1}}\\[0.5cm]
  {\Large Giới thiệu}
  \vfill
\end{frame}

\begin{frame}{Câu hỏi nghiên cứu chính}
  \textbf{Vấn đề cốt lõi:}
  \begin{itemize}
    \item Mạng Neural thường \emph{xấp xỉ tốt} các hàm trơn, nhưng \textbf{thất bại} khi tổng quát hóa trên các phép toán rời rạc.
    \item \textbf{Câu hỏi:} Liệu mạng Neural có thể học thực thi các chỉ dẫn thuật toán mã hóa nhị phân \textbf{một cách chính xác}?
  \end{itemize}

  \vspace{0.3cm}
  \textbf{Phạm vi bài báo:}
  \begin{itemize}
    \item Hoán vị nhị phân, Phép cộng, Phép nhân
    \item Lệnh SBN $\Rightarrow$ \textbf{Turing-complete}
  \end{itemize}
\end{frame}

\begin{frame}{Đóng góp chính của bài báo}
  \textbf{Hai đổi mới quan trọng:}
  \begin{enumerate}
    \item \textbf{Biểu diễn theo mẫu thuật toán:} Dữ liệu huấn luyện dạng "template" - số lượng chỉ tăng \textbf{logarit}.
    \item \textbf{Tính học được chính xác:} Ensemble mạng 2 lớp ($n_h \to \infty$) học \textbf{chính xác} với xác suất cao.
  \end{enumerate}
  
  \vspace{0.3cm}
  $\Rightarrow$ \textbf{Kết quả NTK-based đầu tiên} cho exact learnability.
\end{frame}

% ===== SLIDE CHUYỂN TIẾP SANG PHẦN 2 =====
\begin{frame}{}
  \centering
  \vfill
  {\Huge \textbf{Phần 2}}\\[0.5cm]
  {\Large Chứng minh lý thuyết NTK}
  \vfill
\end{frame}
